\scanpage{153}
\mspage{102}

qu'il n'y a pas de contrexemples avec des catégories \unsure{genre|genres} \deleted{\unsure{$\Delta$}} de la forme $Y=Z_{/F}$ \unsure{$\alpha$} (déduites d'un foncteur $g:X\to Z$ en prenant $F=\unsure{g_!|f_!}(e_X)=\unsure{\pi|\Pi}_0(X/Z)$), où $Z$ est une catégorie telle que $\Delta^{\unsure{n|1}}$\inserted{\unsure{(?)}}$, \Phi,\Psi$, ou

\diagram{\begin{tikzcd}[ampersand replacement=\&,row sep=tiny,column sep=tiny]
\& \cdot \&\\
\cdot \arrow[ur] \arrow[dr] \& \& \cdot \arrow[ul] \arrow[dl]\\
\& \cdot \&
\end{tikzcd}}

(ayant le type d'homotopie du cercle) etc, car les composantes connexes de telles catégories $Y$ sont des $K(\pi,1)$ -- les $\pi_i$ \inserted{supérieurs} $(i\geq 2)$ sont nuls.

Cette difficulté vue, le contrexemple devient évident. Je pars d'un objet \inserted{\unsure{1-connexe|1-connexe ?}} \uline{quelconque} $Y$ de $X$, et je prends $y_0\in\mathrm{ob}\,Y$ et

\[
X=\underline{\mathrm{Ch}}(Y,y_0,-)\xrightarrow{\unsure{t|f}}Y.
\]

\deleted{\unsure{Les}} On sait que $X$ est $\unsure{W|\mathcal W}_\infty$-asph, et que $f$ est $W_\infty$-parfait, donc les inclusions

\[
X_y\unsure{-|\longrightarrow}{}_{y\backslash}X
\]

sont des $W_\infty$-équivalences. L'hyp. a), que les $ {}_{y\backslash}X$ soient 0-connexes revient donc à la $=$ hypothèse sur les fibres\inserted{$X_y$}, ou
% archive sheet number: 153
